\section{Emergent relativity and effective Lorentz symmetry}
\label{sec:emergent-relativity}

The continuum model is written with an external time parameter $t$, so Lorentz invariance is not
postulated at the microscopic level. Nevertheless, in the \emph{linear wave regime} an effective
Lorentz symmetry can emerge as a kinematic symmetry of the dominant dispersion branch.

\subsection{Linear-wave cone and an effective light speed}

Linearizing \Cref{eq:eom} about a homogeneous background yields a multi-branch wave system.
For each polarization branch $b$, the small-amplitude dispersion takes the form
\begin{equation}
  \omega_b^2(\veck)\;\approx\; c_b^2\,|\veck|^2 \quad \text{for } |\veck|\to 0,
  \label{eq:linear-dispersion}
\end{equation}
with branch-dependent characteristic speeds $c_b$ (e.g., transverse and longitudinal speeds in
an isotropic elastic medium). If a single branch dominates the long-distance propagation of
narrowband packets, its wave cone defines an effective ``light speed'' $c:=c_{\mathrm{dom}}$.

\subsection{Effective Minkowski interval}

For the dominant branch, the wavefront condition suggests an effective interval
\begin{equation}
  \mathrm{d}s^2 = c^2\,\mathrm{d}t^2 - \mathrm{d}\mathbf{x}^2 ,
  \label{eq:effective-interval}
\end{equation}
whose null surfaces $\mathrm{d}s^2=0$ coincide with wave cones in the linear regime. The standard
Lorentz transformations are the linear maps that preserve \eqref{eq:effective-interval} and the
orientation of the forward cone:
\begin{align}
t' &= \gamma\left(t - \frac{\mathbf{v}\cdot\mathbf{x}}{c^2}\right), \\
\mathbf{x}' &= \mathbf{x} + (\gamma-1)\frac{(\mathbf{v}\cdot\mathbf{x})\mathbf{v}}{v^2} - \gamma\,\mathbf{v}\,t, \\
\gamma &= (1-v^2/c^2)^{-1/2}.
\end{align}
In this reading, Lorentz symmetry is an \emph{effective} invariance of the linear propagation sector,
analogous to analogue-gravity systems where wave excitations experience an emergent metric
\citep{Barcelo2011AnalogueGravity}.

\subsection{Regime of validity and expected deviations}

The effective Lorentz symmetry is expected to break when:
(i) multiple branches contribute comparably (mode mixing),
(ii) dispersion becomes significant at shorter wavelengths, or
(iii) geometric/material nonlinearity becomes important at larger amplitudes.
These deviations are not defects but concrete targets for validation, because they
predict where the effective field description should fail.
